\subsubsection{射精控制训练与龟头脱敏（自我管理实操）}

行为疗法之所以是早泄的一线治疗，是因为它通过反复练习重建射精反射的阈值与自控感，而非依赖药物。除了经典的挤压法与停-动法（见前文），以下三项自我训练可作为日常补充，坚持数周至数月通常能见到效果。

1. \textbf{射精阈值感知训练（停-动法的单人版）}：
   - 以自慰方式把兴奋程度想象成 0--10 的刻度，训练目标是\textbf{在 7--8 分时主动停下}，待回落到 4--5 分再继续。
   - 每次练习可经历 2--3 次"升—停—降"循环后再射精；每周 3--4 次，持续 6--12 周。
   - 关键在于\textbf{识别"即将高潮"的身体信号}（如会阴收紧、呼吸变浅、龟头敏感陡增），学会在失控前刹车。

2. \textbf{盆底肌（凯格尔）训练}：
   - \textbf{找准肌肉}：排尿时尝试中断尿流所用的肌群即耻骨尾骨肌；也可想象"忍住放屁"的同时轻微上提阴茎根部。
   - \textbf{方法}：收缩 3--5 秒、放松 3--5 秒为一组，每次 10--15 组，每日 2--3 次；熟练后可延长收缩时间。\textbf{注意只收缩盆底肌}，腹部、大腿与臀部尽量放松，勿憋气。
   - \textbf{进阶}：采用"快收快放"（每秒一次）与"慢收慢放"交替；部分男性配合生物反馈训练效果更稳定。
   - \textbf{证据}：多项研究显示盆底肌训练可延长阴道内射精潜伏期（IELT），且无药物副作用，值得作为基础训练坚持。

3. \textbf{龟头脱敏训练（渐进适应）}：
   - \textbf{目的}：降低龟头与冠状沟的过度敏感，提高对刺激的耐受。
   - \textbf{方法}：日常以温水淋浴时翻起包皮、用温水冲洗龟头；在自慰时以手掌或润滑液做\textbf{由轻到重、由慢到快}的渐进刺激，接近高潮即停，逐步延长耐受时间。
   - \textbf{辅助}：必要时可短期使用局部麻醉喷剂/乳膏（见下文药物治疗），但\textbf{务必在性交前清洗干净}，以免麻醉残留影响伴侣或导致自身勃起减退。
   - \textbf{提醒}：脱敏不等于"越刺激越好"，过度粗暴的摩擦可致龟头黏膜损伤，应循序渐进、以不痛为度。

\begin{tcolorbox}[colback=pink!5!white,colframe=red!50!black,title=贴心小叮咛]
自我训练需要耐心：通常以"周"为单位见效，切忌三天打鱼。训练期间与伴侣坦诚沟通，把它当作双方共同的事而非"一个人的功课"。若坚持 8--12 周仍无改善，或伴发勃起困难、射精疼痛、血精、尿路症状，应到泌尿外科/男科进一步评估，排除甲状腺功能异常、前列腺炎等可治病因。
\end{tcolorbox}